The output filename may be an incomplete filename.  In the case of the
\verb|WATCH| point element, this means it may contain one instance of
the string format specification ``\%s'' and one occurence of an
integer format specification (e.g., ``\%ld'').  {\tt elegant} will
replace the former with the rootname (see
\verb|run_setup|) and the latter with the element's occurrence
number.  For example, suppose you had a repetitive lattice defined as
follows:
\begin{verbatim}
W1: WATCH,FILENAME=''%s-%03ld.w1''
Q1: QUAD,L=0.1,K1=1
D: DRIFT,L=1
Q2: QUAD,L=0.1,K1=-1
CELL: LINE=(W1,Q1,D,2*Q2,D,Q1)
BL: LINE=(100*CELL)
\end{verbatim}
The element \verb|W1| appears 100 times.  Each instance will result in
a new file being produced.  Successive instances have names like
``{\em rootname}-001.w1'', ``{\em rootname}-002.w1'', ``{\em
rootname}-003.w1'', and so on up to ``{\em rootname}-100.w1''.  (If
instead of ``\%03ld'' you used ``\%ld'', the names would be ``{\em
rootname}-1.w1'', ``{\em rootname}-2.w1'', etc. up to ``{\em
rootname}-100.w1''.  This is generally not as convenient as the names
don't sort into occurrence order.)

The files can easily be plotted together, as in 
\begin{verbatim}
% sddsplot -column=t,p *-???.w1 -graph=dot -separate 
\end{verbatim}
They may also be combined into a single file, as in
\begin{verbatim}
% sddscombine *-???.w1 all.w1 
\end{verbatim}

In passing, note that if \verb|W1| was defined as
\begin{verbatim}
W1: WATCH,FILENAME=''%s.w1''
\end{verbatim}
or 
\begin{verbatim}
W1: WATCH,FILENAME=''output.w1''
\end{verbatim}
only a single file would be produced, containing output from the last instance
only.

For multi-bunch simulation, a variant of this feature can be used to create multiple files, one for
each bunch. For example,
\begin{verbatim}
W1: WATCH,FILENAME=''%s-%02ld.w1'',BUNCH_SERIES=1
Q1: QUAD,L=0.1,K1=1
D: DRIFT,L=1
Q2: QUAD,L=0.1,K1=-1
CELL: LINE=(Q1,D,2*Q2,D,Q1)
BL: LINE=(100*CELL,48*W1)
\end{verbatim}
would result in the creation of 48 files, one for each of 48 bunches. If more than 48 bunches are present,
they won't be recorded. If fewer are present, the excess files will be empty.

Notes: 
\begin{enumerate}
\item Confusion sometimes occurs about some of the quantities related
  to the {\tt s} coordinate in this file when in parameter mode.  Please
  see Section \ref{sec:longitCoord} above.
\item This element can adversely affect parallel efficiency. Use of the
  \verb|START_PASS|, \verb|END_PASS|, \verb|INTERVAL|,  and \verb|FLUSH_INTERVAL|
  options can help reduce the impact. Also, particle output is the most expensive, by far.
\end{enumerate}
